\clearpage\fancyhead[L]{\small JUNO.}\hypertarget{native-330-lower}{}\label{native-330-lower}
\begin{center}\textbf{JUNO.}\par\vspace{4pt}\hrule\vspace{4pt}\fontsize{10}{12}\selectfont Göttingische gelehrte Anzeigen. Stück 92. 93. S. 913\,..\,914. 1811. Juni 10.\par\vspace{4pt}\hrule\vspace{4pt}\end{center}

\par\indent Die letzten Nachrichten über Juno stehen im 136. Stück unserer Anzeigen von 1808 [S. 305 d. B.]. Dort sind die Beobachtungen mitgeteilt, die Professor \textsc{Gauss} 1808 an unserer Sternwarte an diesem Planeten angestellt hat, außerdem die Bestimmung der vierten Opposition und der achte Satz von Elementen. Die fünfte Opposition, die im Januar 1810 eintrat, wurde nirgends beobachtet. Zum ersten Mal ging die Bestimmung einer Opposition eines der vier neuen Planeten durch ein gemeinsames Versäumnis aller Astronomen verloren. Umso wichtiger war es, zu verhindern, dass auch die \SourceLeaf{331}\emph{sechste} Opposition im April dieses Jahres unbeobachtet blieb. Die wenigen Beobachtungen, die Professor \textsc{Gauss} wegen der äußerst geringen Helligkeit des Planeten nur mit dem Kreismikrometer anstellen konnte, werden daher umso wertvoller sein, da sie allem Anschein nach die einzigen sein werden, die diesmal irgendwo gemacht wurden. Zum Vergleich dienten $\varphi$ Librae und einige andere nahe gelegene Sterne der Histoire céleste. Wenn ihre Positionen künftig genauer bestimmt sind, werden die folgenden Ergebnisse noch eine kleine Korrektur brauchen.
\Needspace{10\baselineskip}
\begin{center}\fontsize{10}{12}\selectfont\setlength{\tabcolsep}{4pt}\renewcommand{\arraystretch}{1.15}\begin{tabular}{c|c|c|c}
1811 & \shortstack{mittlere Zeit\\in Göttingen} & \shortstack{Rektaszension} & \shortstack{Südliche\\Deklination}\\\hline
April 22 & $9^{\text{u}}\;51^{\text{m}}\;35^{\text{s}}$ & $216^\circ\;41'\;50''1$ & $0^\circ\;58'\;16''$\\
24 & $10\;32\;55$ & $216\;17\;58.7$ & $0\;46\;50$\\
25 & $10\;19\;22$ & $216\;6\;3.6$ & $0\;40\;4$\\
26 & $10\;17\;30$ & $215\;54\;7.3$ & $0\;33\;51$\\
\end{tabular}\end{center}

\par\indent Die Deklination vom 22. April ist zweifelhaft, und auch die vom 24. April ist weniger zuverlässig als die beiden folgenden, bei denen der Planet günstiger stand.

\Needspace{8\baselineskip}\par\indent Aus diesen Beobachtungen hat Professor \textsc{Gauss} für die Opposition das folgende Ergebnis gewonnen.
\begin{center}1811 April 24. $19^{\text{u}}\;20^{\text{m}}\;12^{\text{s}}$ mittl. Zeit in Göttingen\\$214^\circ\;8'\;48''3$ wahre heliozentrische Länge\\$12\;55\;2.0$ wahre geozentrische Breite, nördlich\end{center}
\Needspace{20\baselineskip}\par\indent Die starke Abweichung dieser Position von den Elementen, die auf den ersten vier Oppositionen beruhen, zeigt nun auch bei Juno den Einfluss der Störungen, die besonders Jupiter verursacht. Für Juno wird ihre Berechnung einen ebenso ungeheuren Aufwand erfordern wie für Pallas. Es ist nur angemessen, dass diese Arbeit zuerst für die früher entdeckte Pallas erledigt wird. Professor \textsc{Gauss} hat sich daher vorerst damit begnügt, neue elliptische Elemente auf die Oppositionen von 1806, 1807, 1808, 1811 zu stützen. Wir geben sie hier wieder.

\begin{center}\fontsize{10}{12}\selectfont\setlength{\tabcolsep}{4pt}\renewcommand{\arraystretch}{1.15}\begin{tabular}{llr}
\shortstack[l]{Epoche der mittleren Länge, 1811,\\Meridian von Göttingen} & .\;. & $177^\circ\;48'\;1''8$\\
Tägliche mittlere tropische Bewegung & .\;.\;.\;.\;.\;.\;.\;.\;.\;. & $813''2486$\\
Länge des Perihels 1811 & .\;.\;.\;.\;.\;.\;.\;.\;.\;.\;.\;.\;.\;. & $53^\circ\;14'\;32''4$\\
Länge des aufsteigenden Knotens 1811 & .\;.\;.\;.\;.\;.\;.\;.\;.\;. & $171\;9\;13.5$\\
Neigung der Bahn & .\;.\;.\;.\;.\;.\;.\;.\;.\;.\;.\;.\;.\;.\;.\;.\;. & $13\;4\;27.0$\\
Exzentrizität & .\;.\;.\;.\;.\;.\;.\;.\;.\;.\;.\;.\;.\;.\;.\;.\;. & $=\sin14\;44\;9.1$\\
Logarithmus der großen Halbachse & .\;.\;.\;.\;.\;.\;.\;.\;.\;.\;. & $0.4265711.$\\
\end{tabular}\end{center}
